\section{Query construction and adaptive access}\label{sec:C16:adaptive}
A query is a hypothesis about which evidence will help. Query rewriting can add useful bridge concepts, but it can also encode an early mistaken interpretation and narrow the search prematurely. An adaptive agent should preserve uncertainty about its information need and revise the query when returned evidence contradicts its assumptions. The memory access policy is then part of reasoning rather than a fixed frontend.

MRAgent uses a Cue--Tag--Content graph and an active reconstruction loop~\cite{mragent2026}. Tags act as associative intermediates, allowing the model to select relations and then fetch content rather than expanding every neighboring node. The mechanism can use newly discovered cues to change later traversal. Its theoretical separation from restricted passive retrieval depends on the constructed task family and access assumptions; it should not be interpreted as a proof that every active agent outperforms every static retriever. Its published code provides a concrete object for implementation inspection.

JAM separates a hierarchical store retaining raw histories from a runtime Researcher that uses search, open, and browse actions to assemble query-specific context~\cite{jit2026research}. Its training combines verified-trajectory supervision with a subsequent policy-optimization stage. This changes the accounting boundary: less irreversible preprocessing can preserve details, while runtime exploration introduces model calls and latency. Query-conditioned construction does not remove the need for retention, access controls, or calibrated stopping when evidence remains incomplete.

The two mechanisms differ in what is navigated. MRAgent traverses explicit associative relations; JAM navigates a hierarchical workspace with raw pages and summaries. A clean comparison would hold raw evidence, model, maximum accessed content, and total compute fixed while varying the access representation. Without that control, an improvement can arise from more evidence, more reasoning calls, a different training distribution, or a better interface. The architecture's name does not identify the cause.

Adaptive access requires a stopping rule. Let $b$ be the current evidence state and let $a$ be a possible retrieval action. An ideal controller would choose an action when its expected reduction in task loss exceeds its cost. In practice that value is uncertain, and self-assessed sufficiency can be overconfident. A useful implementation distinguishes stopping because evidence is sufficient from stopping because the budget is exhausted. Both can produce an answer, but only the first asserts that the information requirement was met.

A bounded decision model clarifies that stopping rule. Let $\xi\in\Xi$ be a sufficient controlled query state, including the target belief, query-use/history information, source/interface versions and permissions needed to determine available queries, costs, observation kernels and future state. Its target-belief projection is $b_\xi=\Pr(Y\in\cdot\mid\xi)$ for a fixed target $Y$. A posterior over $Y$ alone need not be sufficient: identical target probabilities can coexist with different remaining queries or observation laws. For a fixed finite legitimate action set $\mathcal A$ and loss $\ell:\mathcal A\times\mathcal Y\to[0,L_{\max}]$ with finite $L_{\max}$, set $\mathcal R(b)=\min_{a\in\mathcal A}\mathbb E_b\ell(a,Y)$. A permitted read-only query $q$ has observation kernel $K(\mathrm d o\mid\xi,q)$ and sufficient-state update $T(\xi,q,o)$. Assume Bayes-consistent target posteriors and that query selection uses only $\xi$ and independent randomization, so selection itself supplies no hidden target information. Its one-step information value is
\begin{equation}
 V_q(\xi)=\mathcal R(b_\xi)-
 \mathbb E_{o\mid\xi,q}\mathcal R(b_{T(\xi,q,o)})\geq0.
\end{equation}
The inequality holds because the decision maker may ignore the new observation and retain its former rule. It does not describe a fallible deployed model, which may be distracted or misled by the returned text. It also does not compare different authority sets or changed task objectives.

Let $\mathcal Q(\xi)$ be a finite available query family, $B\geq0$ a finite resource budget, $c(\xi,q)\geq0$ the declared query cost, and $\lambda\geq0$ a finite loss-scale multiplier. Write $\mathcal Q_B(\xi)=\{q\in\mathcal Q(\xi):c(\xi,q)\leq B\}$. An explicit remaining-query horizon $n\in\mathbb N_0$ gives the finite-horizon controller
\begin{equation}
\begin{aligned}
 W_0(\xi,B)&=\mathcal R(b_\xi),\\
 D_n(\xi,B,q)&=\lambda c(\xi,q)+\mathbb E_{o\mid\xi,q}
 W_{n-1}(T(\xi,q,o),B-c(\xi,q)),\\
 W_n(\xi,B)&=\min\!\left\{\mathcal R(b_\xi),
 \min_{q\in\mathcal Q_B(\xi)}D_n(\xi,B,q)\right\},\quad n\geq1.
\end{aligned}
\end{equation}
With no feasible query the inner minimum is $+\infty$. Unavailable results remain observations with incurred cost. The explicit horizon terminates the recursion even when a permitted query has zero modeled cost. A budget alone bounds query count only under an additional positive lower cost bound. A myopic comparison $V_q(\xi)>\lambda c(\xi,q)$ can miss complementary multi-query evidence; the recursion exposes that limitation without claiming that a correct sufficient state or observation model is available in a production agent. Log budget or horizon exhaustion separately from a verified evidence-sufficiency decision.

The proposed diagnostic logs each information need, query, selected relation or page, returned evidence, and reason for stopping. It then replays failed cases with an oracle bridge cue or a missing page supplied. Repair under that intervention supports an access-policy bottleneck. Failure despite the evidence shifts attention to interpretation or action selection. This analysis is more informative than counting how many search calls an agent made.

The released code makes several boundaries more concrete. In the inspected MRAgent \path{Agent.answer_question} path, question-key extraction precedes graph-relation evaluation; temporal questions can take a specialized branch, and non-temporal evidence undergoes coarse selection, identifier-based deduplication, finer selection, and key-tag expansion before the tool-enabled answer stage~\cite{repo_mragent2026}. These operations mean that a graph match is not identical to final evidence exposure. Deduplication by identifier prefix can reduce repeated material while also changing how much supporting detail survives. A reproduction should log both pre-selection candidates and the final answer input.

The repository linked by JAM includes a research implementation whose \texttt{ResearchAgent.research} refreshes retrievers, loads memory, plans a search, accumulates a temporary result, and reflects on sufficiency~\cite{repo_jam2026}. It records iterations and stops either when the reflection says enough or when the configured iteration bound ends. The returned integrated memory alone does not distinguish those stopping causes; the raw iteration record is needed to reconstruct them. The inspected March 2026 repository revision precedes the September JAM report. Its research and application variants therefore do not, by themselves, establish the deployed weights, training path, or every runtime interface of the paper's trained configuration. Public code and paper mechanism must be mapped at that narrower boundary.

MATM's inspected \texttt{UnifiedRerankerOrchestrator.rerank} supports staged reranking, normalizes candidates into one schema, and retains stage-level error metadata~\cite{repo_matm2026}. In the inspected error path, a failed stage can preserve its upstream candidate order instead of producing a new ranking. Thus a run configured for learned reranking need not actually apply that ranking on every query. Reporting configured mode without stage application and error counters can misattribute a result to the reranker. These are source observations at pinned revisions, not measured error frequencies or reproduced benchmark outcomes.

